\documentclass[10pt]{article}
\usepackage[margin=0.65in]{geometry}
\usepackage[T1]{fontenc}
\usepackage{lmodern,microtype,parskip,enumitem,xcolor,hyperref}
\definecolor{accent}{HTML}{0F6D67}
\hypersetup{colorlinks=true,urlcolor=accent}
\setlist[itemize]{leftmargin=1.2em,itemsep=1pt,topsep=1pt}
\newcommand{\cvsec}[1]{\vspace{5pt}{\color{accent}\rule{\linewidth}{.6pt}}\vspace{-5pt}\section*{#1}\vspace{-5pt}}
\newcommand{\e}[2]{\textbf{#1}\hfill #2\\}
\begin{document}
\begin{center}{\LARGE\bfseries Hector Bahamonde}\\ Senior Researcher in Political Science\\ University of Turku, Finland\\ \href{mailto:hector.bahamonde@utu.fi}{hector.bahamonde@utu.fi}\\ CV date: 14 August 2026\end{center}
\cvsec{1. Personal details}
Political scientist specialising in comparative politics, political economy, inequality, causal inference, and experimental methods. Languages: English and Spanish (native). Advanced software: R, Stata, Python, oTree, and \LaTeX.
\cvsec{2. Degrees}
\e{Title of Docent in Political Science, University of Turku}{2024}\e{Ph.D. in Political Science, Rutgers University, USA}{2017}\e{B.A. in Political Science, Catholic University of Chile}{2009}
\cvsec{3. Current employment}
\e{Senior Researcher, INVEST, University of Turku}{Sep. 2021--present}\e{Project Researcher, Åbo Akademi University}{Apr. 2026--present}\e{Research Director, INVEST-Experimental Research Hub}{2022--2023}
\cvsec{4. Previous work experience}
\e{Assistant Professor, O'Higgins University, Chile}{2018--2021}\e{Postdoctoral Fellow, Tulane University}{2017--2018}\e{Teaching Assistant, Rutgers University}{2014--2017}
\cvsec{5. Research funding and grants}
Excellence Fellowship, Rutgers University (2012--2017); Small Grant Fund for Research on Latin America (2013); Jerome M. Clubb Scholarship, ICPSR (2013); Pre-Dissertation Award, Rutgers (2016); Experimental Workshop Award (2015).
\cvsec{6. Research output}
Research interests: democratic backsliding, authoritarian values, clientelism, populism, state formation, inequality, comparative political economy, political development, causal inference, and experimental methods.
\textbf{Selected peer-reviewed publications}\begin{itemize}\item Bahamonde and Sarpila (2024), Physical appearance and elections, \emph{Political Psychology}.\item Bahamonde and Canales (2022), Electoral Risk and Vote Buying, \emph{Electoral Studies}.\item Bahamonde (2022), Still for Sale, \emph{Acta Politica}.\item Bahamonde and Trasberg (2021), Inclusive Institutions, Unequal Outcomes, \emph{European Journal of Political Economy}.\item Modrego, Canales, and Bahamonde (2020), Employment Effects of COVID-19, \emph{Regional Science Policy \& Practice}.\item Bahamonde (2018), Aiming Right at You, \emph{Journal of Politics in Latin America}.\end{itemize}
\textbf{Current manuscripts and projects}\begin{itemize}\item First Impressions under Scarcity (work in progress).\item Losers, Delegation, and the Ideological Valence of Expertise: Evidence from Finland.\item Losers' Consent and Anti-Systemic Protest: Experimental Evidence from Chile and Estonia.\item Vote-Selling and Vote-Buying (under review).\end{itemize}
\cvsec{7. Methods, software, data, and tools}
Advanced R, Stata, Python, oTree, and \LaTeX user; experienced in survey, laboratory, natural, and conjoint experiments, causal inference, panel/time-series, maximum likelihood, and Bayesian analysis.
\cvsec{8. Research supervision and leadership}
PhD supervision at the University of Turku: Davit Totadze and Rugaiia Bairamova (2025--present). Research Director, INVEST-Experimental Research Hub (2022--2023). Coordinator, INVESThub Brown Bag and Workshop series (2022--present).
\cvsec{9. Teaching merits}
University of Turku: Quantitative Methods; Experimental Methods in Social Sciences. O'Higgins University: Quantitative Methods I--II, Political Science I--II, Social Sciences and Epistemology. Tulane and Rutgers: comparative politics, quantitative methods, American government, and computing.
\cvsec{10. Awards, honours, and training}
Rutgers Excellence Fellowship; Pre-Dissertation Award; Experimental Workshop Award; Jerome M. Clubb Scholarship. Certified training in experimental political science, causal inference, Bayesian analysis, advanced econometrics, time series, and behavioral economics.
\cvsec{11. Memberships and peer review}
Editorial Board Member, \emph{Social Sciences} (Springer Nature), 2020--present. Referee for the \emph{American Political Science Review}, \emph{Journal of Politics}, \emph{Comparative Political Studies}, \emph{Political Behavior}, and other journals. Panel organiser, chair, or discussant at APSA, ECPR, EPSA, MPSA, LASA, and FPSA.
\cvsec{12. Invited lectures and conferences}
Invited seminars at the University of Turku, Åbo Akademi, INVEST, the Centre for Experimental Social Sciences, and SOMA's Appearance \& Inequalities Seminar. Conference participation, 2010--2026: APSA, EPSA, ECPR, MPSA, LASA, SPSA, FPSA, and NOLAN.
\cvsec{13. Scientific and societal impact}
Research communicated through the \emph{New York Times}, CNN Chile, and Chilean regional media, addressing democratic inequality, elections, vote buying, employment, and poverty.
\cvsec{14. Other relevant merits}
Book review in the \emph{Journal of Iberian and Latin American Research}; member and organiser of international scholarly networks.
\cvsec{15. References}
Robert Kaufman (Rutgers); Daniel Kelemen (Georgetown); Paul Poast (Chicago); Juan Pablo Luna (McGill); Ludovico Feoli (Tulane); Ross Baker (Rutgers).
\end{document}
